\documentclass[12pt]{article}
\everymath{\displaystyle}
\usepackage{unicode-math}
\usepackage{pgfplots}
\pgfplotsset{compat=1.18}
\usepackage{amsmath}



\usepackage[margin=1in]{geometry}
\usepackage{tcolorbox}
\usepackage{graphicx}
\usepackage{tikz}
\usepackage{amssymb}
\usepackage{enumitem}
\usepackage{titlesec}
\usepackage{fancyhdr}
\usepackage{pifont}
\usepackage{xcolor}
\newcounter{questionnum}
\setcounter{questionnum}{0}
\usepackage{eso-pic} % for page background

% --------- 主题颜色设置 ---------
\definecolor{novelblue}{RGB}{70,60,150}
\definecolor{headerbg}{RGB}{240,240,255}

\pagestyle{fancy}
\fancyhf{}
\renewcommand{\headrulewidth}{0.4pt}
\renewcommand{\headrule}{\hbox to\headwidth{\color{novelblue}\leaders\hrule height 0.4pt\hfill}}

% --------- 页眉背景色条 ---------
\newcommand{\headbackground}{
  \AddToShipoutPictureBG*{
    \begin{tikzpicture}[remember picture, overlay]
      \fill[blue!5!purple!10] (current page.north west) rectangle ([yshift=-60pt]current page.north east);
    \end{tikzpicture}
  }
}

% --------- 页眉设置 ---------
\pagestyle{fancy}
\fancyhf{}
\renewcommand{\headrulewidth}{0pt}
\fancyhead[L]{\color{novelblue}\textbf{\textbf{Novel Prep} }}
\fancyhead[C]{\color{novelblue}\textbf {SAT Preparation}}
\fancyhead[R]{\color{novelblue}\textbf{Jack Zeng}}


\newtcolorbox{SATCard}[1][]{%
  enhanced,
  colback=white,
  colframe=novelblue,
  coltitle=white,
  boxrule=0.6pt,
  arc=4pt,
  boxsep=5pt,
  left=8pt, right=8pt, top=10pt, bottom=10pt,
  sharp corners=south,
  drop shadow south east,
  #1
}

% --------- 顶部 Mark for Review 栏 ---------
\newcommand{\markforreview}{
  \stepcounter{questionnum} % 自动递增题号
  \begin{tikzpicture}[scale=1.1]
    % 黑色题号
    \fill[black] (0,0) rectangle (0.8,0.6);
    \node[white, font=\bfseries\large] at (0.4,0.3) {\thequestionnum};

    % 灰底主栏
    \fill[gray!10] (0.8,0) rectangle (14.5,0.6);

    % Mark for Review
    \node[anchor=west, font=\small] at (1.2,0.3) {\ding{72} \hspace{2pt}Mark for Review};

    % ABC 按钮区域
    \draw[thick, rounded corners=3pt] (13.5,0.12) rectangle (14.5,0.48);
    \node[font=\bfseries\normalsize] at (14.0,0.3) {ABC};
  \end{tikzpicture}


  \newcommand{\choiceboxcircled}[2]{%
  \begin{tcolorbox}[
    colback=white,
    colframe=black,
    boxrule=0.5pt,
    rounded corners=6pt,
    left=6pt, right=6pt, top=6pt, bottom=6pt,
    width=\linewidth
  ]
    \textbf{\large\textcircled{\textbf{#1}}} \ #2
  \end{tcolorbox}
}
}

% --------- 选项框样式 ---------
\newcommand{\choicebox}[2]{%
  \begin{tcolorbox}[
    colback=white, 
    colframe=novelblue,
    boxrule=0.5pt, 
    rounded corners=4pt, 
    left=6pt, right=6pt, top=6pt, bottom=6pt,
    width=\linewidth
  ]
  \textbf{#1.} \ #2
  \end{tcolorbox}
}


\begin{document}
\section*{Math Hard Question 11}
\noindent\markforreview
\vspace{1em}

In the given equation, \(a\) and \(b\) are positive constants.  
The sum of the solutions to the given equation is \(k(6a + b)\), where \(k\) is constant.  
What is the value of \(k\)?  
\[
24x^2 - (12a + 2b)x + ab = 0
\]

\choicebox{A}{\(\frac{1}{12}\)}  
\choicebox{B}{\(-\frac{1}{12}\)}  
\choicebox{C}{12}  
\choicebox{D}{-12}

\vspace{2em}
\noindent\markforreview
\vspace{1em}

The positive number \(a\) is 2136\% of the sum of the positive numbers \(b\) and \(c\),  
and \(b\) is 89\% of \(c\). What percent of \(b\) is \(a\)?

\choicebox{A}{4536\%}  
\choicebox{B}{2400\%}  
\choicebox{C}{4037\%}  
\choicebox{D}{2225\%}

\vspace{2em}
\newpage
\noindent\markforreview
\vspace{1em}

The graphs in the \(xy\)-plane of the following quadratic equations each have  
\(x\)-intercepts of \(-2\) and 4. The graph of which equation has its vertex  
**farthest from the x-axis**?

\choicebox{A}{\(y = -7(x + 2)(x - 4)\)}  
\choicebox{B}{\(y = \frac{1}{10}(x + 2)(x - 4)\)}  
\choicebox{C}{\(y = -\frac{1}{2}(x + 2)(x - 4)\)}  
\choicebox{D}{\(y = 5(x + 2)(x - 4)\)}

\vspace{2em}
\noindent\markforreview
\vspace{1em}

In one hour, the distance, \(d(s)\), in kilometers that a ferry can travel up and down a river  
flowing with a constant speed \(s\), in meters per second, is:  
\[
d(s) = 10.7 - 1.2s^2
\]  
where \(s\) and \(d(s)\) are positive.  

Which of the following equivalent expressions for \(d(s)\) contains the speed  
of the river in meters per second, as a constant or coefficient, for which  
the ferry can travel a distance of 8 kilometers in one hour?

\choicebox{A}{\(8 - 1.2(s - 1.5)(s + 1.5)\)}  
\choicebox{B}{\(8 + 1.2(2.25 - s^2)\)}  
\choicebox{C}{\(8(1 - 0.15s^2) + 2.7\)}  
\choicebox{D}{\(11.9 - 1.2(s - 1)^2 - 2.4s\)}

\vspace{2em}
\newpage
\noindent\markforreview
\vspace{1em}

A minor league hockey team has been collecting ticket sales data over the past year.  
At a current price of \$25 per ticket, an average of 4000 seats are purchased.  
They predict that for each \$1 increase in ticket price, 100 fewer tickets will be sold.  

Which of the following functions best models the amount of money, \(y\),  
that the hockey team expects to collect from ticket sales based on an \(x\)-dollar increase in ticket price?

\choicebox{A}{\(y = (25 + x)(4000 - 100x)\)}  
\choicebox{B}{\(y = (25 - x)(4000 + 100x)\)}  
\choicebox{C}{\(y = x(4000 - 100x)\)}  
\choicebox{D}{\(y = 4000(25 + x)\)}

\vspace{2em}
\noindent\markforreview
\vspace{1em}

In the given expression:  
\[
34z^{14} + bz^7 + 70,
\]  
\(b\) is a positive integer. If \(qz^7 + r\) is a factor of the expression, where \(q\) and \(r\) are positive integers,  
what is the greatest possible value of \(b\)?



\vspace{2em}
\newpage

\noindent\markforreview
\vspace{1em}

Point \(N\) lies on a unit circle in the \(xy\)-plane and has coordinates \((1, 0)\).  
Point \(O\) is the center of the circle and has coordinates \((0, 0)\).  
Point \(M\) also lies on the unit circle, and the measure of angle \(\angle NOM\) is \(\frac{266\pi}{3}\) radians.  
If the coordinates of point \(M\) are \((a, b)\), where \(a\) and \(b\) are constants, what is the value of \(a\)?

\choicebox{A}{\(\frac{1}{2}\)}  
\choicebox{B}{\(\frac{\sqrt{3}}{2}\)}  
\choicebox{C}{\(-\frac{1}{2}\)}  
\choicebox{D}{\(-\frac{\sqrt{3}}{2}\)}

\vspace{2em}
\noindent\markforreview
\vspace{1em}

For each of two rectangles, the ratio of the rectangle’s length to its width is \(10 : 5\).  
The width of rectangle \(X\) is 16.5 times the width of rectangle \(Y\).  
How does the length of rectangle \(X\) compare to the length of rectangle \(Y\)?

\choicebox{A}{The length of rectangle \(X\) is 0.5 times the length of rectangle \(Y\)}  
\choicebox{B}{The length of rectangle \(X\) is 2 times the length of rectangle \(Y\)}  
\choicebox{C}{The length of rectangle \(X\) is 16.5 times the length of rectangle \(Y\)}  
\choicebox{D}{The length of rectangle \(X\) is 33 times the length of rectangle \(Y\)}

\vspace{2em}

\newpage
\noindent\markforreview
\vspace{1em}

In the \(xy\)-plane, a circle has center \(C\) with coordinates \((h, k)\).  
Points \(A\) and \(B\) lie on the circle. Point \(A\) has coordinates \((h + 1, k + \sqrt{102})\),  
and \(\angle ACB\) is a right angle. What is the length of \(\overline{AB}\)?

\choicebox{A}{\(\sqrt{206}\)}  
\choicebox{B}{\(2\sqrt{102}\)}  
\choicebox{C}{\(103\sqrt{2}\)}  
\choicebox{D}{\(103\sqrt{3}\)}

\vspace{2em}

\noindent\markforreview
\vspace{1em}

\[
y = a * (\frac{6}{b})^{x}-c

\]

How many times does the graph of the given equation in the \(xy\)-plane cross the \(x\)-axis,  
where \(a\), \(b\), and \(c\) are positive constants such that \(a > 6\) and \(b > c\)?

\choicebox{A}{Zero}  
\choicebox{B}{One}  
\choicebox{C}{Two}  
\choicebox{D}{Three}

\vspace{2em}
\newpage
\noindent\markforreview
\vspace{1em}

Two lines intersect at exactly one point, forming two acute angles and two obtuse angles.  
The measure of one of these angles is \(9x - 110^\circ\).  
Which of the following could **NOT** be the sum of the measures of any two of these angles?

\choicebox{A}{\(-18x + 220^\circ\)}  
\choicebox{B}{\(-18x + 580^\circ\)}  
\choicebox{C}{\(18x - 220^\circ\)}  
\choicebox{D}{\(180^\circ\)}

\vspace{2em}
\noindent\markforreview
\vspace{1em}

The function \(f\) is defined by:  
\[
f(x) = a(3.1^x + 3.1^b),
\]  
where \(a\) and \(b\) are integer constants and \(0 < a < b\).  
The functions \(g\) and \(h\) are equivalent to function \(f\), where \(k\) and \(m\) are constants.  

Which of the following equations displays the **\(y\)-coordinate of the \(y\)-intercept** of the graph  
of \(y = f(x)\) in the \(xy\)-plane as a constant or coefficient?

\begin{itemize}
  \item[I.] \(g(x) = a(3.1^x + k)\)
  \item[II.] \(h(x) = a(3.1^x) + m\)
\end{itemize}

\choicebox{A}{I only}  
\choicebox{B}{II only}  
\choicebox{C}{I and II only}  
\choicebox{D}{Neither I nor II}

\vspace{2em}









\end{document}